\begin{defn}[Adjacency Matrix]
    An \emph{adjacency matrix} for a graph is a two dimensional array \(A\) in which \(a_{ij}\) is the number of edges that we may traverse from vertex \(i\) to vertex \(j\).
\end{defn}

% \begin{practice}[sidebyside]{Euler Circuits}

% \end{practice}

% \begin{practice}[sidebyside]{Euler Trails}
% \begin{center}
%     \begin{tikzpicture}
%         \graph[
%             tree layout,
%             grow'=right,
%             edges={edges,rounded corners},
%             nodes={vertex},
%             sibling distance = 1.5cm,
%             level distance = 1.5cm
%         ]{
%             A--{C,B,A};
%             D--{E};
%             D --[bend angle=45,bend left] A;
%             E--B;
%             B -- D;
%             C -- B;
%         };
%     \end{tikzpicture}
% \end{center}
% \end{practice}

\begin{expository}[sidebyside]{Small Adjacency Example}
\begin{center}
    \begin{tikzpicture}
        \graph[
            tree layout,
            grow'=right,
            edges={edges,rounded corners},
            nodes={vertex},
            sibling distance = 3cm,
            level distance = 3cm
        ]{
            A--[bend left] B,
            B--[bend left] A,
            A-- C,
            C-- B,
            % C -- {B,D};
        };
    \end{tikzpicture}
\end{center}
\tcblower
\[
\begin{blockarray}{cccc}
    & A & B & C \\
\begin{block}{c(ccc)}
    A & \phantom{\rule{1cm}{1cm}} & \phantom{\rule{1cm}{1cm}} & \phantom{\rule{1cm}{1cm}} \\
    B & \phantom{\rule{1cm}{1cm}} & \phantom{\rule{1cm}{1cm}} & \phantom{\rule{1cm}{1cm}} \\
    C & \phantom{\rule{1cm}{1cm}} & \phantom{\rule{1cm}{1cm}} & \phantom{\rule{1cm}{1cm}} \\
\end{block}
\end{blockarray}
 \]

\end{expository}

\begin{practice}[sidebyside]{Another Small-ish Example}
\begin{center}
    \begin{tikzpicture}
        \graph[
            tree layout,
            grow'=right,
            edges={edges,rounded corners},
            nodes={vertex},
            sibling distance = 1.75cm,
            level distance = 2cm
        ]{
            A--{C,B,A},
            B -- D,
            C -- {B,D},
            C --[loop above] C,
        };
    \end{tikzpicture}
\end{center}
\tcblower

\[
\begin{blockarray}{ccccc}
    & A & B & C & D \\
\begin{block}{c(cccc)}
    A & \phantom{\rule{1cm}{1cm}} & \phantom{\rule{1cm}{1cm}} & \phantom{\rule{1cm}{1cm}} & \phantom{\rule{1cm}{1cm}} \\
    B & \phantom{\rule{1cm}{1cm}} & \phantom{\rule{1cm}{1cm}} & \phantom{\rule{1cm}{1cm}} & \phantom{\rule{1cm}{1cm}} \\
    C & \phantom{\rule{1cm}{1cm}} & \phantom{\rule{1cm}{1cm}} & \phantom{\rule{1cm}{1cm}} & \phantom{\rule{1cm}{1cm}} \\
    D & \phantom{\rule{1cm}{1cm}} & \phantom{\rule{1cm}{1cm}} & \phantom{\rule{1cm}{1cm}} & \phantom{\rule{1cm}{1cm}} \\
\end{block}
\end{blockarray}
 \]

\end{practice}

\begin{practice}[sidebyside]{Directed Graph Example}
\begin{center}
    \begin{tikzpicture}
        \graph[
            layered layout,
            grow'=right,
            edges={edges,rounded corners,every loop/.style={->}},
            nodes={vertex},
            sibling distance = 1.75cm,
            level distance = 2cm,
        ]{
            A->{C,B,A},
            B -> C,
            C ->[loop above] C,
            D[desired at={(1,0)}] ->[bend right] E[desired at={(1,1)}],
            E ->[bend right] D,
        };
    \end{tikzpicture}
\end{center}
\tcblower

\[
\begin{blockarray}{cccccc}
    & A & B & C & D & E\\
\begin{block}{c(ccccc)}
    A & \phantom{\rule{1cm}{1cm}} & \phantom{\rule{1cm}{1cm}} & \phantom{\rule{1cm}{1cm}} & \phantom{\rule{1cm}{1cm}} & \phantom{\rule{1cm}{1cm}}\\
    B & \phantom{\rule{1cm}{1cm}} & \phantom{\rule{1cm}{1cm}} & \phantom{\rule{1cm}{1cm}} & \phantom{\rule{1cm}{1cm}} & \phantom{\rule{1cm}{1cm}}\\
    C & \phantom{\rule{1cm}{1cm}} & \phantom{\rule{1cm}{1cm}} & \phantom{\rule{1cm}{1cm}} & \phantom{\rule{1cm}{1cm}} & \phantom{\rule{1cm}{1cm}}\\
    D & \phantom{\rule{1cm}{1cm}} & \phantom{\rule{1cm}{1cm}} & \phantom{\rule{1cm}{1cm}} & \phantom{\rule{1cm}{1cm}} & \phantom{\rule{1cm}{1cm}}\\
    E & \phantom{\rule{1cm}{1cm}} & \phantom{\rule{1cm}{1cm}} & \phantom{\rule{1cm}{1cm}} & \phantom{\rule{1cm}{1cm}} & \phantom{\rule{1cm}{1cm}} \\
\end{block}
\end{blockarray}
 \]

\end{practice}

\clearpage

\begin{expository}{Matrix Arithmetic}
We add matrices term by term:
\begin{center}
    \begin{tikzpicture}
    % Product
        \node at (0,0){\(
            \begin{pmatrix*}[r]
                1 & 2 \\
                0 & 3 
            \end{pmatrix*}+
            \begin{pmatrix*}[r]
                1 & 7\\
                0 & 2
            \end{pmatrix*}=
            \begin{pmatrix*}[r]
                1+1 & 2+7 \\
                0+0 & 3+2 \\
            \end{pmatrix*}
        \)};
    % Labels
        \draw[secondaryaccent,ultra thick] (-2.7,0.3) circle (7pt) node (a){};
        \draw[secondaryaccent,ultra thick] (-1,0.3) circle (7pt) node (b){};
        \draw[->,secondaryaccent,ultra thick,shorten >=3pt,shorten <=3pt] 
            (a) to[bend angle=60,bend left] node[pos=0.5,above] {+} (b);
    \end{tikzpicture}
\end{center}
We multiply matrices combining rows with columns:
\begin{center}
    \begin{tikzpicture}
    % Product
        \node at (0,0){\(
            \begin{pmatrix*}[r]
                1 & 2 \\
                0 & 3 
            \end{pmatrix*}\ 
            \begin{pmatrix*}[r]
                1 & 0 & 7\\
                0 & -1 & 2
            \end{pmatrix*}=
            \begin{pmatrix*}[r]
                1\cdot 1+2\cdot 0 & 1\cdot 0+2\cdot -1 & 1\cdot 7+2\cdot 2\\
                0\cdot 1+3\cdot 0 & 0\cdot 0+3\cdot -1 & 0\cdot 7+3\cdot 2\\
            \end{pmatrix*}
        \)};
    % Labels
        \draw[->,secondaryaccent,ultra thick] (-5.7,0.8) -- +(1,0);
        \draw[->,secondaryaccent,ultra thick] (-4.35,0.7) -- +(0,-1.3);
    \end{tikzpicture}
\end{center}
\end{expository}

\begin{practice}{Matrix Arithmetic}
Find the sum, difference, and product of the following matrices where possible.
\begin{equation*}
    A=\begin{pmatrix*}[r]
        2 & 7\\ 1 & 4
    \end{pmatrix*},\ 
    B=\begin{pmatrix*}[r]
        4 & -7\\ -1 & 2
    \end{pmatrix*},\ 
    C=\begin{pmatrix*}[r]
        2 & 0 & 1\\ 1 & 0 & -1
    \end{pmatrix*},\ 
    I=\begin{pmatrix*}[r]
        1 & 0\\ 0 & 1
    \end{pmatrix*}
\end{equation*}
\begin{enumerate}
    \setlength\itemsep{2.4\baselineskip}
    \item \(A+B=\)
    \item \(B-A=\)
    \item \(AB=\)
    \item \(AI=\)
    \item \(AC=\)
    \item \(CA=\)
    \item \(C+B=\)\vspace{2.4\baselineskip}
\end{enumerate}

\end{practice}

\clearpage

\begin{expository}[sidebyside,lefthand width=0.2\textwidth]{Counting Walks}
\begin{center}
    \begin{tikzpicture}[scale=0.6,transform shape]
        \graph[
            tree layout,
            grow'=right,
            edges={edges,rounded corners},
            nodes={vertex},
            sibling distance = 3cm,
            level distance = 3cm
        ]{
            A--[bend left] B,
            B--[bend left] A,
            A--C,
            C--B,
            % C -- {B,D};
        };
    \end{tikzpicture}
\end{center}
\[
    A=\begin{pmatrix*}[r]
        0 & 2 & 1\\
        2 & 0 & 1\\
        1 & 1 & 0
    \end{pmatrix*}
\]

\tcblower
\begin{enumerate}
    \item Walks of length 1:\vspace{5\baselineskip}
    \item Walks of length 2:\vspace{5\baselineskip}
    \item Walks of length 3:\vspace{5\baselineskip}
\end{enumerate}

\end{expository}

\begin{practice}[sidebyside,lefthand width=0.2\textwidth]{Counting Walks}
\begin{center}
    \begin{tikzpicture}[scale=0.6,transform shape]
        \graph[
            tree layout,
            grow'=right,
            edges={edges,rounded corners,every loop/.style={->},shorten >=5pt},
            nodes={vertex},
            sibling distance = 3cm,
            level distance = 3cm
        ]{
            A->[bend left] B,
            B<-[bend left] A,
            A--[loop above] A,
            A-> C,
            C-> B,
        };
    \end{tikzpicture}
\end{center}
\[
    A=\begin{pmatrix*}[r]
        1 & 2 & 1\\
        0 & 0 & 0\\
        0 & 1 & 0
    \end{pmatrix*}
\]

\tcblower
\begin{enumerate}
    \item Walks of length 1:\vspace{5\baselineskip}
    \item Walks of length 2:\vspace{5\baselineskip}
    \item Walks of length 3:\vspace{5\baselineskip}
\end{enumerate}

\end{practice}

\clearpage

\begin{expository}{Concatenating (``Adding'') Graphs}
\begin{center}
    \begin{tikzpicture}[
        every node/.style={vertex},
        every edge/.style={draw,edges}
    ]
    \begin{scope}
    % nodes
        \node (a) at (0,0) {\(A\)};
        \node (b) [below=of a] {\(B\)};
        \node (c) [right=of b] {\(C\)};
    % edges
        \path (a) edge (b);
        \path (b) edge (c);
    \end{scope}
        \node[shade=none,text width=2cm] (P) [above right=of c,yshift=-0.75cm] {\Huge \(+\)};
    \begin{scope}[xshift=6cm]
    % nodes
        \node (a) at (0,0) {\(A\)};
        \node (b) [below=of a] {\(B\)};
        \node (c) [right=of b] {\(C\)};
    % edges    
        \path (a) edge (c);
    \end{scope}
        \node[shade=none,text width=2cm] (E) [above right=of c,yshift=-0.75cm] {\Huge \(=\)};
    \begin{scope}[xshift=12cm]
    % nodes
        \node (a) at (0,0) {\(A\)};
        \node (b) [below=of a] {\(B\)};
        \node (c) [right=of b] {\(C\)};
    % edges    
        \path (a) edge (c);
        \path (a) edge (b);
        \path (b) edge (c);
    \end{scope}
    \end{tikzpicture}
\end{center}

\tcblower

\begin{equation*}
    \begin{pmatrix*}[r]
        0&1&0\\
        1&0&1\\
        0&1&0\\
    \end{pmatrix*}
        +
    \begin{pmatrix*}[r]
        0&0&1\\
        0&0&0\\
        1&0&0\\
    \end{pmatrix*}
        =
    \begin{pmatrix*}[r]
        0&1&1\\
        1&0&1\\
        1&1&0\\
    \end{pmatrix*}
\end{equation*}
\end{expository}

\begin{practice}{Concatenating Graphs}
\begin{center}
    \begin{tikzpicture}[
        every node/.style={vertex},
        every edge/.style={draw,edges}
    ]
    \begin{scope}
    % nodes
        \node (a) at (0,0) {\(A\)};
        \node (b) [below=of a] {\(B\)};
        \node (c) [right=of b] {\(C\)};
    % edges
        \path (a) edge (b) edge[bend left] (c);
        \path (b) edge (c);
        \path (c) edge[loop right] (c);
    \end{scope}
        \node[shade=none,text width=2cm] (P) [above right=of c,yshift=-0.75cm] {\Huge \(+\)};
    \begin{scope}[xshift=6cm]
    % nodes
        \node (a) at (0,0) {\(A\)};
        \node (b) [below=of a] {\(B\)};
        \node (c) [right=of b] {\(C\)};
    % edges    
        \path (a) edge[bend right] (c);
        \path (a) edge[loop above] (a);
    \end{scope}
        \node[shade=none,text width=2cm] (E) [above right=of c,yshift=-0.75cm] {\Huge \(=\)};
    \begin{scope}[xshift=12cm]
    % nodes
        \node (a) at (0,0) {\(A\)};
        \node (b) [below=of a] {\(B\)};
        \node (c) [right=of b] {\(C\)};
    % edges    
        % \path (a) edge (c);
        % \path (a) edge (b);
        % \path (b) edge (c);
    \end{scope}
    \end{tikzpicture}
\end{center}

\tcblower

\vspace{20\baselineskip}


\end{practice}

\clearpage

\begin{expository}{Directed Graphs, Adjacency Matrices, and Websites}
\begin{center}
    \begin{tikzpicture}[
        every node/.style={inner sep=2pt,align=center},
        every edge/.style={edges,draw,every loop/.style={->}},
        scale=1,transform shape
    ]
    % Sites
        \node[text width=1in,] (a) at (0,0) {Alice's Restaurant};
        \node[text width=0.5in,] (b) [below right=of a] {Bob's Burger};
        \node[text width=1in,] (f) [below left=of a] {Frying Dutchman};
        \node (e) [below right =of f] {Elzar's};
    % Links
        \path[->] (a) edge[bend left] (b)
            edge[bend right] (f);
        \path[->] (b) edge[bend left] (a);
        \path[->] (e) edge (a)
            edge[loop right] (e);
        \path[->] (f) edge (a)
            edge[bend angle=10,bend right] (b)
            edge[bend right] (e)
            edge[loop left] (f);
    \end{tikzpicture}
\end{center}
\begin{equation*}
    L=
    \begin{blockarray}{ccccc}
        & A & B & E & F \\
    \begin{block}{c(cccc)}
        A & 0 & 1/2 & 0 & 1/2 \\
        B & 1 & 0 & 0 & 0 \\
        E & 1/2 & 0 & 1/2 & 0 \\
        F & 1/4 & 1/4 & 1/4 & 1/4 \\
    \end{block}
    \end{blockarray}
    =
    \begin{pmatrix*}[r]
        0 & 0.5 & 0 & 0.5 \\
        1 & 0 & 0 & 0 \\
        0.5 & 0 & 0.5 & 0\\
        0.25 & 0.25 & 0.25 & 0.25
    \end{pmatrix*}
\end{equation*}

\begin{minted}{python}
import numpy as np
from numpy import linalg as LA

# matrix of links/probabilities
L = np.array([[0,0.5,0,0.5],[1,0,0,0],[0.5,0,0.5,0],[0.25,0.25,0.25,0.25]])

# Initial position vector
P=np.array([1,0,0,0])

distance = 1
i = 0
while distance>0.001 and i<100:
    # Multiply by L to represent making a random click
    temp_P=P@L
    # Compare the position vectors
    distance = LA.norm(temp_P-P)
    # Update
    P=temp_P
    i+=1
print(P)
\end{minted}
\begin{equation*}
    P=\begin{bmatrix}
        0.37519702 & 0.24988937 & 0.12502424 & 0.24988937
    \end{bmatrix}    
\end{equation*}
Starting at the website for Alice's Restaurant and then randomly clicking around there is a 37\% chance we get back to Alice's, 25\% chance of ending up at Bob's or the Dutchman, and only  13\% chance of ending up at Elzar's website. Thus, we rank Alice's site highest and Elzar's lowest.

    % \(L\) is the link matrix, \(L^n\) is the naive approach
    % \(dL+(1-d)J\) is the link matrix combined with another matrix to add random jumps, \(J\) is all ones
    % Look in classes/code/PageRank.ipynb for example from Imperial College example.
    
\end{expository}

